\section{Analysis and Design for Topological Robustness}
\label{sec:topological}

A structural property the primitive lacks, a sensitivity analysis, then a closed-form-plus-numerical
design section.

\subsection{The Connectivity Blind Spot}
\label{subsec:blindspot}

This subsection formally states a property DTAF1 \textbf{lacks}: pointwise EDT tolerance-matching
is not ``connectivity-injective.'' Concretely (proof sketch; full statement in
Appendix~\ref{app:proofs}):

\begin{quote}
For a road mask $G_c$ of width $w$ and tolerance $d_c$, there exist predicted masks $P_c$ obtained
from $G_c$ by deleting an arbitrarily large fraction of pixels --- up to and including masks in
which $G_c$'s connected skeleton is fully disconnected into many isolated fragments --- such that
$\mathrm{DTAF1}_c(P_c, G_c) = 1.0$, provided the surviving pixels remain within $d_c$ of every GT
pixel.
\end{quote}

Section~\ref{subsec:motivation-results}'s road-breakage sweep is the empirical witness: DTAF1
stays at exactly $1.000$ through 90\% random dropout on the synthetic scene.

\subsection{APLS as Connectivity Recovery, and Its Sensitivity to Fragmentation}
\label{subsec:apls-recovery}

APLS (Section~\ref{subsec:apls}) recovers exactly the signal DTAF1 lacks: its score depends on
shortest-path length \emph{between} control points, not merely on point-to-point proximity, so it
is sensitive to whether a path exists at all. Framing the dropout fraction as the ``noise'' axis,
the empirical curve has a two-regime character: a near-flat, low-sensitivity regime while the
skeleton is still mostly connected, followed by a sharp departure once fragmentation actually
occurs (Table~\ref{tab:breakage-topo}).

\begin{table}
  \caption{Road-breakage sweep: DTAF1 vs.\ CBHM vs.\ APLS vs.\ DTAF1-Topo.}
  \label{tab:breakage-topo}
  \begin{tabular}{@{}lcccc@{}}
    \toprule
    Dropout & \texttt{dtaf1} & \texttt{cbhm} & \texttt{road\_apls} & \texttt{dtaf1\_topo} \\
    \midrule
    0.0 & 1.000 & 1.000 & 1.000 & 1.000 \\
    0.1 & 1.000 & 0.965 & 0.898 & 0.946 \\
    0.2 & 1.000 & 0.919 & 0.883 & 0.938 \\
    0.3 & 1.000 & 0.909 & \textbf{0.247} & \textbf{0.397} \\
    0.5 & 1.000 & 0.805 & 0.029 & 0.057 \\
    0.7 & 1.000 & 0.606 & 0.005 & 0.010 \\
    1.0 & 0.500 & 0.000 & 0.000 & 0.000 \\
    \bottomrule
  \end{tabular}
\end{table}

The transition is sharp between 0.2 and 0.3: \texttt{road\_apls} falls from $0.883$ to $0.247$ in
that one step --- the point at which the road skeleton's largest connected fragment actually
detaches from enough control-point pairs that shortest paths start disappearing rather than merely
lengthening. \texttt{dtaf1} shows no analogous transition anywhere in the swept range.

\subsection{Harmonic-Blending Design for DTAF1-Topo}
\label{subsec:dtaf1topo}

\begin{lstlisting}[language=Python, caption={Blending primitive, \texttt{metrics/dtaf1\_topo.py:29-31,74-87}}]
def _harmonic_mean(a, b):
    denom = a + b
    return float(2 * a * b / denom) if denom > 0 else 0.0

# per class in linear_classes:
entry["blended"] = _harmonic_mean(base["f1"], a)  # a = apls
# non-linear classes keep entry["blended"] = base["f1"]
\end{lstlisting}

\subsubsection{Closed-Form Boundary Behavior}

The harmonic mean of $F1$ and $\mathrm{APLS}$ has a simple, exactly-analyzable shape at its
boundaries:
\begin{itemize}
  \item $\mathrm{APLS}=0, F1=1$ (the DTAF1 blind-spot regime, worst case): $\mathrm{blended}=0$. A
    perfectly tolerance-matched but fully disconnected road scores exactly zero once blended,
    regardless of how high $F1$ is.
  \item $\mathrm{APLS}=1, F1=1$: $\mathrm{blended}=1$, recovering DTAF1's own perfect score
    unchanged.
  \item For fixed $F1=1$, $\mathrm{blended}(\mathrm{APLS}) = 2\mathrm{APLS}/(1+\mathrm{APLS})$ ---
    concave and always $\le \mathrm{APLS}$, meaning the harmonic blend never \emph{rewards}
    connectivity beyond what APLS itself reports; it can only ever penalize a
    topologically-correct-looking $F1$ down toward the (generally lower) APLS score, never the
    reverse. This asymmetry is intentional: it is the mechanism by which DTAF1-Topo closes the
    blind spot rather than merely averaging two independent opinions.
\end{itemize}

\subsubsection{Numerical Validation via the Road-Breakage Sweep}

Unlike the closed-form boundary analysis above, the \emph{practical} question --- where along the
dropout axis does the blend actually diverge from plain DTAF1? --- has no closed form and is
answered numerically. Table~\ref{tab:breakage-topo} already gives the answer directly:
\texttt{dtaf1\_topo} tracks \texttt{road\_apls} closely (both driven to near-zero by fraction
0.5--0.7) while \texttt{dtaf1} stays at $1.000$ throughout --- this \emph{is} the fix, cross-checked
against the assertions in \texttt{tests/test\_metrics\_sanity.py::TestRoadBreakageRegression}.
\textbf{Validated on synthetic data only} --- see Section~\ref{sec:discussion} for the real-data
caveat.

\medskip
\noindent\textbf{Bridge to Section~\ref{sec:method}.} Every property analyzed so far --- offset
tolerance (Section~\ref{subsec:dtaf1}), width insensitivity (Section~\ref{subsec:cldice}),
connectivity/breakage sensitivity (this section), instance-level point matching
(Section~\ref{subsec:pointf1}) --- was developed purely as an \emph{evaluation} concern: a
property a \emph{finished} prediction is scored against, after training is already done.
Section~\ref{sec:method} asks the natural next question: if these are the geometric properties
that actually distinguish a good extraction from a bad one, why train a model with a loss that has
no notion of any of them? The three tractable properties (tolerance, width-insensitivity via
topology, point-instance-like localization) are exactly the three \texttt{losses/} implements as
differentiable training terms; the one property with no tractable differentiable relaxation (graph
connectivity, this section) is exactly the one \texttt{losses/} leaves out, by design, not
oversight.
